% GENERATED FILE. Edit canonical lessons, then run npm run generate:books.
% canonical-source-hash: fnv1a64:8077fe371146bc46
% canonical-lessons: ML-C165-vrttiyulla, ML-C165-madhuramulla, ML-C165-kaypulla, ML-C165-erivulla, ML-C165-puliyulla

\chapter{Clean, Sweet, Bitter, Spicy, Sour}
\label{ch:ml-clean-sweet-bitter-spicy-sour}
\hlchaptermodality{165}

\begin{chapteropening}
\textbf{By the end of this chapter:} \emph{I can say clean, sweet, bitter, spicy and sour.}

Complete the last lesson of chapter 165: I can say clean, sweet, bitter, spicy and sour.
\end{chapteropening}

\section[\texorpdfstring{\ml{വൃത്തിയുള്ള} (vṛttiyuḷḷa)}{വൃത്തിയുള്ള (vṛttiyuḷḷa)}]{\ml{വൃത്തിയുള്ള} (vṛttiyuḷḷa) — clean}

\label{lesson:ML-C165-vrttiyulla}



\begin{quote}
Before the new one: say the Malayalam for cold, then the Malayalam for long.

In \emph{muppatŭ}, \textquotedblleft{}thirty,\textquotedblright{} is the three in its usual shape? (No: an older, shorter form.)
\end{quote}

\subsection*{You'll want to know: \ml{വൃത്തിയുള്ള}}
\textbf{\ml{വൃത്തിയുള്ള}} — \emph{vṛttiyuḷḷa} — \textquotedblleft{}clean\textquotedblright{}.

\begin{grammarlens}[title={What you've built}]
One more word for this chapter.
\end{grammarlens}

\subsection*{Your turn}
\begin{itemize}
\raggedright
  \item \emph{Say it:} \emph{vṛttiyuḷḷa}
  \item \emph{Say it:} \emph{vṛttiyuḷḷa}, once more
  \item \emph{Say it:} say \emph{vṛttiyuḷḷa}
  \item \emph{Recall:} say the Malayalam for cold, then the Malayalam for long, then say \emph{vṛttiyuḷḷa} again
\end{itemize}

\begin{tcolorbox}[breakable,colback=teal!4,colframe=teal!35!black,title={Before you move on}]
What does \ml{വൃത്തിയുള്ള} mean? (\textquotedblleft{}Clean\textquotedblright{}.) Say it once more.
\end{tcolorbox}
\section[\texorpdfstring{\ml{മധുരമുള്ള} (madhuramuḷḷa)}{മധുരമുള്ള (madhuramuḷḷa)}]{\ml{മധുരമുള്ള} (madhuramuḷḷa) — sweet}

\label{lesson:ML-C165-madhuramulla}



\begin{quote}
Before the new one: say the Malayalam for long, then the Malayalam for clean.
\end{quote}

\subsection*{You'll want to know: \ml{മധുരമുള്ള}}
\textbf{\ml{മധുരമുള്ള}} — \emph{madhuramuḷḷa} — \textquotedblleft{}sweet\textquotedblright{}.

\begin{grammarlens}[title={What you've built}]
One more word for this chapter.
\end{grammarlens}

\subsection*{Your turn}
\begin{itemize}
\raggedright
  \item \emph{Say it:} \emph{madhuramuḷḷa}
  \item \emph{Say it:} \emph{madhuramuḷḷa}, once more
  \item \emph{Say it:} say \emph{madhuramuḷḷa}
  \item \emph{Recall:} say the Malayalam for long, then the Malayalam for clean, then say \emph{madhuramuḷḷa} again
\end{itemize}

\begin{tcolorbox}[breakable,colback=teal!4,colframe=teal!35!black,title={Before you move on}]
What does \ml{മധുരമുള്ള} mean? (\textquotedblleft{}Sweet\textquotedblright{}.) Say it once more.
\end{tcolorbox}
\section[\texorpdfstring{\ml{കയ്പുള്ള} (kaypuḷḷa)}{കയ്പുള്ള (kaypuḷḷa)}]{\ml{കയ്പുള്ള} (kaypuḷḷa) — bitter}

\label{lesson:ML-C165-kaypulla}



\begin{quote}
Before the new one: say the Malayalam for clean, then the Malayalam for sweet.

Say \textquotedblleft{}This is not beautiful.\textquotedblright{} (\emph{Itŭ sundaramalla}.)
\end{quote}

\subsection*{You'll want to know: \ml{കയ്പുള്ള}}
\textbf{\ml{കയ്പുള്ള}} — \emph{kaypuḷḷa} — \textquotedblleft{}bitter\textquotedblright{}.

\begin{grammarlens}[title={What you've built}]
One more word for this chapter.
\end{grammarlens}

\subsection*{Your turn}
\begin{itemize}
\raggedright
  \item \emph{Say it:} \emph{kaypuḷḷa}
  \item \emph{Say it:} \emph{kaypuḷḷa}, once more
  \item \emph{Say it:} say \emph{kaypuḷḷa}
  \item \emph{Recall:} say the Malayalam for clean, then the Malayalam for sweet, then say \emph{kaypuḷḷa} again
\end{itemize}

\begin{tcolorbox}[breakable,colback=teal!4,colframe=teal!35!black,title={Before you move on}]
What does \ml{കയ്പുള്ള} mean? (\textquotedblleft{}Bitter\textquotedblright{}.) Say it once more.
\end{tcolorbox}
\section[\texorpdfstring{\ml{എരിവുള്ള} (erivuḷḷa)}{എരിവുള്ള (erivuḷḷa)}]{\ml{എരിവുള്ള} (erivuḷḷa) — spicy}

\label{lesson:ML-C165-erivulla}



\begin{quote}
Before the new one: say the Malayalam for sweet, then the Malayalam for bitter.
\end{quote}

\subsection*{You'll want to know: \ml{എരിവുള്ള}}
\textbf{\ml{എരിവുള്ള}} — \emph{erivuḷḷa} — \textquotedblleft{}spicy\textquotedblright{}.

\begin{grammarlens}[title={What you've built}]
One more word for this chapter.
\end{grammarlens}

\subsection*{Your turn}
\begin{itemize}
\raggedright
  \item \emph{Say it:} \emph{erivuḷḷa}
  \item \emph{Say it:} \emph{erivuḷḷa}, once more
  \item \emph{Say it:} say \emph{erivuḷḷa}
  \item \emph{Recall:} say the Malayalam for sweet, then the Malayalam for bitter, then say \emph{erivuḷḷa} again
\end{itemize}

\begin{tcolorbox}[breakable,colback=teal!4,colframe=teal!35!black,title={Before you move on}]
What does \ml{എരിവുള്ള} mean? (\textquotedblleft{}Spicy\textquotedblright{}.) Say it once more.
\end{tcolorbox}
\section[\texorpdfstring{\ml{പുളിയുള്ള} (puḷiyuḷḷa)}{പുളിയുള്ള (puḷiyuḷḷa)}]{\ml{പുളിയുള്ള} (puḷiyuḷḷa) — sour}

\label{lesson:ML-C165-puliyulla}



\begin{quote}
Before the new one: say the Malayalam for bitter, then the Malayalam for spicy.

Say \textquotedblleft{}neither a teacher nor a doctor.\textquotedblright{} (\emph{Adhyāpakanum vaidyanum alla}.)
\end{quote}

\subsection*{You'll want to know: \ml{പുളിയുള്ള}}
\textbf{\ml{പുളിയുള്ള}} — \emph{puḷiyuḷḷa} — \textquotedblleft{}sour\textquotedblright{}.

\begin{grammarlens}[title={What you've built}]
One more word for this chapter.
\end{grammarlens}

\subsection*{Your turn}
\begin{itemize}
\raggedright
  \item \emph{Say it:} \emph{puḷiyuḷḷa}
  \item \emph{Say it:} \emph{puḷiyuḷḷa}, once more
  \item \emph{Say it:} say \emph{puḷiyuḷḷa}
  \item \emph{Recall:} say the Malayalam for bitter, then the Malayalam for spicy, then say \emph{puḷiyuḷḷa} again
\end{itemize}

\begin{tcolorbox}[breakable,colback=teal!4,colframe=teal!35!black,title={Before you move on}]
What does \ml{പുളിയുള്ള} mean? (\textquotedblleft{}Sour\textquotedblright{}.) Say it once more.
\end{tcolorbox}
